\documentclass[11pt,a4paper]{article}
% preamble.tex — estilo común de todos los apuntes Froth.
% Cada apunte hace % preamble.tex — estilo común de todos los apuntes Froth.
% Cada apunte hace % preamble.tex — estilo común de todos los apuntes Froth.
% Cada apunte hace \input{preamble} justo después de \documentclass.
\usepackage[margin=2.4cm]{geometry}
\usepackage{xcolor}
\definecolor{litblue}{RGB}{31,79,121}
\definecolor{litgreen}{RGB}{27,94,32}
\definecolor{litorange}{RGB}{230,81,0}
\usepackage{parskip}
\usepackage{titlesec}
\titleformat{\section}{\Large\bfseries\color{litblue}}{\thesection}{0.6em}{}
\titleformat{\subsection}{\large\bfseries\color{litblue}}{\thesubsection}{0.6em}{}
\usepackage{tcolorbox}
\usepackage{fancyhdr}
\pagestyle{fancy}
\fancyhf{}
\fancyhead[L]{\small\color{litblue}Froth — Apuntes de ML}
\fancyhead[R]{\small\thepage}
\renewcommand{\headrulewidth}{0.4pt}
\usepackage[colorlinks=true,linkcolor=litblue,urlcolor=litblue]{hyperref}

% Cabecera de cada apunte: \cabecera{numero}{titulo}
\newcommand{\cabecera}[2]{%
  \begin{tcolorbox}[colback=litblue,coltext=white,colframe=litblue,arc=2mm]
    {\Large\bfseries Apunte #1 — #2}\\[3pt]
    {\small Proyecto Froth · aprender Machine Learning construyendo}
  \end{tcolorbox}\vspace{0.8em}}

% Cajas temáticas
\newtcolorbox{concepto}[1]{colback=blue!4,colframe=litblue,title={Concepto: #1},arc=1.5mm}
\newtcolorbox{practica}{colback=green!5,colframe=litgreen,title={Pruébalo tú (teclado en mano)},arc=1.5mm}
\newtcolorbox{ojo}{colback=orange!8,colframe=litorange,title={Ojo / error típico},arc=1.5mm}
\newtcolorbox{resumen}{colback=gray!8,colframe=gray!60,title={En una frase},arc=1.5mm}

\newcommand{\codigo}[1]{\texttt{#1}}
 justo después de \documentclass.
\usepackage[margin=2.4cm]{geometry}
\usepackage{xcolor}
\definecolor{litblue}{RGB}{31,79,121}
\definecolor{litgreen}{RGB}{27,94,32}
\definecolor{litorange}{RGB}{230,81,0}
\usepackage{parskip}
\usepackage{titlesec}
\titleformat{\section}{\Large\bfseries\color{litblue}}{\thesection}{0.6em}{}
\titleformat{\subsection}{\large\bfseries\color{litblue}}{\thesubsection}{0.6em}{}
\usepackage{tcolorbox}
\usepackage{fancyhdr}
\pagestyle{fancy}
\fancyhf{}
\fancyhead[L]{\small\color{litblue}Froth — Apuntes de ML}
\fancyhead[R]{\small\thepage}
\renewcommand{\headrulewidth}{0.4pt}
\usepackage[colorlinks=true,linkcolor=litblue,urlcolor=litblue]{hyperref}

% Cabecera de cada apunte: \cabecera{numero}{titulo}
\newcommand{\cabecera}[2]{%
  \begin{tcolorbox}[colback=litblue,coltext=white,colframe=litblue,arc=2mm]
    {\Large\bfseries Apunte #1 — #2}\\[3pt]
    {\small Proyecto Froth · aprender Machine Learning construyendo}
  \end{tcolorbox}\vspace{0.8em}}

% Cajas temáticas
\newtcolorbox{concepto}[1]{colback=blue!4,colframe=litblue,title={Concepto: #1},arc=1.5mm}
\newtcolorbox{practica}{colback=green!5,colframe=litgreen,title={Pruébalo tú (teclado en mano)},arc=1.5mm}
\newtcolorbox{ojo}{colback=orange!8,colframe=litorange,title={Ojo / error típico},arc=1.5mm}
\newtcolorbox{resumen}{colback=gray!8,colframe=gray!60,title={En una frase},arc=1.5mm}

\newcommand{\codigo}[1]{\texttt{#1}}
 justo después de \documentclass.
\usepackage[margin=2.4cm]{geometry}
\usepackage{xcolor}
\definecolor{litblue}{RGB}{31,79,121}
\definecolor{litgreen}{RGB}{27,94,32}
\definecolor{litorange}{RGB}{230,81,0}
\usepackage{parskip}
\usepackage{titlesec}
\titleformat{\section}{\Large\bfseries\color{litblue}}{\thesection}{0.6em}{}
\titleformat{\subsection}{\large\bfseries\color{litblue}}{\thesubsection}{0.6em}{}
\usepackage{tcolorbox}
\usepackage{fancyhdr}
\pagestyle{fancy}
\fancyhf{}
\fancyhead[L]{\small\color{litblue}Froth — Apuntes de ML}
\fancyhead[R]{\small\thepage}
\renewcommand{\headrulewidth}{0.4pt}
\usepackage[colorlinks=true,linkcolor=litblue,urlcolor=litblue]{hyperref}

% Cabecera de cada apunte: \cabecera{numero}{titulo}
\newcommand{\cabecera}[2]{%
  \begin{tcolorbox}[colback=litblue,coltext=white,colframe=litblue,arc=2mm]
    {\Large\bfseries Apunte #1 — #2}\\[3pt]
    {\small Proyecto Froth · aprender Machine Learning construyendo}
  \end{tcolorbox}\vspace{0.8em}}

% Cajas temáticas
\newtcolorbox{concepto}[1]{colback=blue!4,colframe=litblue,title={Concepto: #1},arc=1.5mm}
\newtcolorbox{practica}{colback=green!5,colframe=litgreen,title={Pruébalo tú (teclado en mano)},arc=1.5mm}
\newtcolorbox{ojo}{colback=orange!8,colframe=litorange,title={Ojo / error típico},arc=1.5mm}
\newtcolorbox{resumen}{colback=gray!8,colframe=gray!60,title={En una frase},arc=1.5mm}

\newcommand{\codigo}[1]{\texttt{#1}}

\usepackage{tikz}
\begin{document}
\cabecera{12}{Reducción de dimensiones: de 768 a 2 sin perder el vecindario}

\section{El problema}

Cada paper es un punto en \textbf{768 dimensiones} (Apuntes 05--07). Tu pantalla tiene
\textbf{2}. Para dibujar el mapa hay que ``aplastar'' — pero un mal aplastado encima puntos
que no tienen nada que ver, y el mapa engaña. La pregunta de este apunte:
\emph{¿cómo se aplasta conservando quién está cerca de quién?}

\section{La intuición clásica: PCA y las sombras}

\begin{concepto}{PCA (análisis de componentes principales)}
Piensa en la \textbf{sombra} de tu mano en la pared: un objeto 3D proyectado en 2D. Si
orientas bien la mano, la sombra conserva la forma (se ven los dedos); si la orientas mal,
todo se encima. \textbf{PCA busca el ``ángulo de luz'' que menos información pierde}: las
direcciones en las que los datos más varían. Es la técnica clásica, rápida y muy usada.
\end{concepto}

\begin{center}
\begin{tikzpicture}[scale=0.9]
  \fill[litblue!60] (0.3,2.1) circle (3pt);
  \fill[litblue!60] (0.8,2.5) circle (3pt);
  \fill[litblue!60] (0.5,2.9) circle (3pt);
  \fill[litorange!80] (2.6,0.6) circle (3pt);
  \fill[litorange!80] (3.1,1.0) circle (3pt);
  \fill[litorange!80] (2.9,0.3) circle (3pt);
  \draw[->,gray,thick] (-0.5,0) -- (4,0) node[below left] {\small proyección buena};
  \draw[litblue] (0.3,0.08) -- (0.9,0.08);
  \draw[litorange] (2.5,0.08) -- (3.2,0.08);
  \node[gray] at (1.7,-0.75) {\small los dos grupos siguen separados};

  \begin{scope}[xshift=7cm]
    \fill[litblue!60] (0.5,2.3) circle (3pt);
    \fill[litblue!60] (1.0,2.7) circle (3pt);
    \fill[litblue!60] (0.7,3.1) circle (3pt);
    \fill[litorange!80] (0.4,0.6) circle (3pt);
    \fill[litorange!80] (0.9,1.0) circle (3pt);
    \fill[litorange!80] (0.7,0.3) circle (3pt);
    \draw[->,gray,thick] (-0.5,0) -- (4,0) node[below left] {\small proyección mala};
    \draw[litblue] (0.45,0.10) -- (1.05,0.10);
    \draw[litorange] (0.35,0.06) -- (0.95,0.06);
    \node[gray] at (1.7,-0.75) {\small ¡los grupos se enciman!};
  \end{scope}
\end{tikzpicture}
\end{center}

Limitación: PCA solo hace proyecciones \emph{rectas} (lineales). Si los datos forman nubes
curvas o enredadas —como suele pasar con texto—, hasta el mejor ángulo encima cosas.

\section{La moderna: UMAP y la mudanza de vecindarios}

\begin{concepto}{UMAP}
UMAP no proyecta sombras. Trabaja en dos tiempos:
\begin{enumerate}
  \item En 768D, apunta para cada paper \textbf{quiénes son sus vecinos más cercanos}.
  \item Coloca los puntos en 2D doblando y estirando lo que haga falta, con una sola
        obsesión: \textbf{que cada punto conserve a sus vecinos al lado}.
\end{enumerate}
Como mudar una ciudad a un mapa más pequeño garantizando que cada casa siga junto a las
mismas casas de antes, aunque las distancias largas cambien. Por eso UMAP separa nubes que
PCA enciman: no está limitado a proyecciones rectas.
\end{concepto}

\section{La letra pequeña: todo mapa 2D miente un poco}

\begin{ojo}
Al aplastar 768 dimensiones a 2 \textbf{siempre} se pierde algo. En un mapa UMAP:
\begin{itemize}
  \item \textbf{Confía en}: quién está pegado a quién; las nubes/grupos que se forman.
  \item \textbf{NO confíes en}: las distancias largas (``este cluster está al doble de
        distancia que aquel'') ni en el tamaño aparente de las nubes.
  \item Los ejes X e Y \textbf{no significan nada} (no son ``más flotación'' o ``más
        litio''); son solo el resultado del acomodo. Por eso los mapas UMAP se publican
        sin números en los ejes.
\end{itemize}
Saber esto te pone por delante de mucha gente que interpreta estos mapas de más.
\end{ojo}

\section{Cómo lo usaremos (lo que viene)}

\begin{enumerate}
  \item Instalar \codigo{umap-learn} (si Python 3.14 lo permite; hay plan B con PCA/t-SNE
        de scikit-learn — cambia la pieza, no el concepto).
  \item Correr UMAP sobre la matriz $126 \times 768$ $\rightarrow$ tabla de 126 puntos (x, y).
  \item Dibujar el primer scatter: tu primer vistazo al mapa (Apunte 13).
  \item Después, agrupar las nubes en subtemas con HDBSCAN (Apunte 14).
\end{enumerate}

\begin{resumen}
PCA aplasta buscando el mejor ``ángulo de sombra'' (rápido pero recto); UMAP reacomoda
preservando los vecindarios (ideal para texto). En el mapa resultante: confía en los
vecinos y las nubes, desconfía de distancias largas, y los ejes no significan nada.
\end{resumen}

\end{document}
